\chapter{Finite Support and Selection}

Chapter 6 began the finite calculus by reading the shape of a measurement: when a reading is
stationary, when its first-order response vanishes, and how that vanishing becomes a detector.
But stationarity quietly assumed something it did not supply --- a field of candidates for the
reading to be stationary over. Before a certificate can declare a residual quiet, it must name
the field in which the residual may be tested. This chapter supplies that field, and supplies
it finitely. It gives the apparatus a finite support to read, coordinate names for the
support, local equations the named support can carry, a measured way to spend the support, and
a tower of refinements that squeezes the residual below any bound. It is the chapter in which
finitude stops being an apology and becomes expressive enough to carry equations.

The first move is to put support before order. A finite support is the held field of
candidates the apparatus permits itself to inspect --- not the whole continuum, not a hidden
reservoir behind a finite notation, but exactly the company a certificate brings to the table.
Support says which candidates count as available; order says only how the apparatus visits
them. The two are different structures, and keeping support ahead of order is what stops the
apparatus from smuggling membership into the act of reading. A candidate is read because it
already stood in the field, not because the trace happened to touch it.

The second move is to name the field without importing analysis. Once the apparatus holds a
finite support, it can place the supported candidates among coordinate names --- a finite
addressing discipline that says where a candidate stands and how adjacent names face one
another, without turning the support into a continuum or promising that every limiting
question has already been answered. Coordinates give the field names; they do not give it an
analysis it has not earned.

The third move is to let the named field carry equations. With a finite support and
coordinate names, the apparatus can write local equations for the middle of a trace --- the
cubic middle-node equations that make the supported record articulate. These are not global
theorems about all paths; they are local equations for the names the support has made
available, a small algebraic scene in which residual quiet can be inspected. Finitude has
become expressive: the field carries equations rather than merely listing names.

The fourth move is to spend the support by a measured selector. A finite support may hold
several admitted candidates, and the apparatus needs a disciplined way to spend the current
support rather than an arbitrary one. Heartbeat selection supplies it --- a measured choice
among admitted candidates by a carried finite cost, allowing ties to remain visible and
passing the remaining support forward. Selection settles a local question, which admitted
candidate this finite event reads now, without pretending to settle the deeper question of
what the residual pressure means across many such readings.

The fifth move is to squeeze the residual along a tower. A refining family of finite supports
can carry receipts forward, compare residual magnitudes across stages, and drive a terminal
residual below a bound that vanishes along the admitted sequence. The squeeze is real
discipline: it lets the apparatus say a residual is small, and smaller along this admitted
tower, without pretending one finite stage has become the whole continuum. Control of residual
magnitude is earned; interpretation of the residual is not yet.

Together these five moves give the apparatus a finite field that is held, named,
equation-bearing, measurably spent, and squeezed toward zero. The chapter does not escape
finitude; it makes finitude carry the weight an analysis would otherwise carry. And it closes,
in its bridge, by drawing the line its own success exposes: the apparatus can now control a
residual along a tower, but control is not yet interpretation, and the question of what single
quantity reads the remaining pressure is the question the next chapter must take up.

\section{Support Before Order}

A completed certificate can close a residual only after the apparatus has
named the field in which the residual may be tested. The earlier chapter
explained what stationarity says once a middle, a path, and a finite
certificate already stand in hand. It separated the direct residual from the
first-order residual, and it showed how the same first-order quiet can appear
as balance, spline closure, or a tridiagonal certificate residual when the
finite identifications have already been carried. That account still leaves a
prior question. Before the certificate can declare quiet, which alternatives
does it read?

Support answers that question. A finite support is the held field of candidates
that the apparatus permits itself to inspect. It does not pretend to contain
the whole continuum. It does not hide an infinite reservoir behind a finite
notation. It does not apologize for leaving uncarried possibilities outside
the present act. It says something sharper: for this certificate, these are
the candidates in scope. The certificate may test them, order them, consume
them, and report what remains. A claim becomes tangible when the apparatus can
hold this finite field, and it becomes fungible when another compatible event
can carry the same field without reconstructing the entire origin of each
candidate.

Write the support as
\[
  K = \{b_1,\ldots,b_n\}.
\]
The notation marks a finite field of candidate middles, candidate variations,
or candidate slips, depending on the surrounding problem. The letter \(K\)
does not name an interval, a region of analysis, or a hidden space of possible
answers. It names the finite company the certificate actually brings to the
table. A candidate belongs to the present field precisely when
\[
  b \in K.
\]
This membership statement carries the first discipline of Chapter 7. The
apparatus may read \(b\) because \(b\) stands inside the carried support. If a
candidate does not stand inside \(K\), this certificate has not yet earned the
right to speak about it.

Support comes before order. The support tells us which candidates count as
available; order tells us how the apparatus visits them. These are different
pieces of structure. The same support can admit several execution orders,
because membership does not say which candidate comes first, which comes next,
or which one closes the trace. If three candidates stand in \(K\), the
apparatus may inspect them in one sequence or another while preserving the same
membership field. Changing the order changes the trace. It need not change the
support.

This distinction keeps the finite doctrine honest. If order created support,
then the apparatus would smuggle membership into the act of reading. A
candidate would count only because the trace happened to touch it. Measurement
requires the stricter order: the field must already stand in hand before the
trace consumes from it. The trace does not produce permission by moving. It
uses permission already carried by the support. The event that reads a
candidate therefore has two checks. First, the candidate belongs to \(K\).
Second, the execution chooses that candidate at this step.

The basic event has the form
\[
  K \xrightarrow{\,b\,} K',
  \qquad
  b \in K,
  \qquad
  K' \subseteq K,
  \qquad
  |K'| < |K|.
\]
Read this display from left to right. The apparatus begins with a finite
support \(K\). It consumes the supported candidate \(b\). It carries forward a
remaining support \(K'\). The remaining support stays inside the original
field, and its size decreases. Nothing here says that the apparatus has solved
for every conceivable candidate. The display says that one finite event has
selected one candidate that already belonged to the field, and that the event
has made the field smaller.

This decrease matters. A finite trace needs a clock it can spend. Without the
inequality \(|K'| < |K|\), consumption would not differ from rehearsal. The
apparatus could visit a candidate and leave the same unresolved field behind.
With the inequality, each consumption event has cost. The field loses one
available candidate, or at least becomes strictly shorter in the sense the
certificate tracks. The trace cannot consume forever while claiming to remain
inside one finite support. Each step spends part of the field it inherited.

The subset condition matters just as much. The remaining candidates are not
new arrivals. The apparatus does not consume \(b\) and then admit a fresh
candidate from outside the field to replace it. Every candidate left in \(K'\)
already stood in \(K\). The trace may reorder, postpone, and eventually
consume supported candidates, but it may not silently enlarge the field it was
given. Enlargement would require a new act of support, not a hidden side
effect of execution.

These two clauses let the apparatus separate membership from sequence. The
support says what may be read. The execution order says which supported
candidate the apparatus reads now. A finite certificate can then record the
trace without confusing the field with the itinerary. This matters whenever a
later calculation presents an analytic-looking result. The calculation may
look as if it has swept across an open landscape, but the finite certificate
has not done that. It has read a carried field in an order.

The residual inherits the same discipline. Once the trace closes, it may
report terminal quiet only for the candidates whose membership the support
carried. The public statement has the form
\[
  r_{\mathrm{term}}(b')=0
  \qquad \text{for } b' \in K.
\]
The quantifier lives inside the support. The prime on \(b'\) reminds us that
the statement concerns any candidate still considered as a supported
variation, not only the candidate consumed at one particular step. The
terminal residual vanishes on supported variations. It does not thereby speak
for unsupported variations, uncarried middles, or a background continuum.

When the stationarity language from the previous chapter enters, the same
support-local sentence becomes
\[
  r_{\mathrm{spline}}(b')=0
  \qquad \text{for } b' \in K.
\]
This display does not add a new universe of candidates. It translates terminal
quiet into the spline residual language already earned by the carried finite
identifications. The support condition remains visible. The apparatus may say
that the spline residual vanishes among the carried candidates. It may not drop
the phrase "for \(b' \in K\)" and pretend that a finite trace has surveyed
every possible variation.

The word "possible" needs care here. Measurement never forbids the operator
from opening a larger support, nor from comparing two supports, nor from
forming a later limiting promise. It only refuses to let one finite
certificate claim more than it carries. A larger support is a new field. A new
field needs its own membership. A limit-promise needs its own finite grammar
for how one field follows another. None of those later acts can erase the
local accounting of the present trace.

Nor does support-local speech weaken the residual. It tells the apparatus where
the residual has a name. A residual without a support can sound larger than the
apparatus that reports it. It can suggest a voice floating above every
candidate whether or not the candidate has entered the finite record. The
support prevents that drift. It ties the residual to the candidates that the
apparatus has made available for inspection. When the certificate says quiet,
it says quiet there: inside the field the trace carried, among the candidates
whose membership the support made tangible, and along the order by which the
trace spent that field. A later certificate may compare fields. It may explain
how one support refines another, or how a family of supports keeps a stable
promise. But this certificate earns no such family by implication. It earns
the finite sentence written on its carried field.

This discipline also gives failure a place to stand. If a candidate outside
\(K\) matters, the apparatus must say so by opening another support or by
showing how the present support should expand. The objection cannot hide
inside the closed residual, because the residual never claimed that outside
candidate. It must appear as a new support question. Measurement therefore
turns overclaim into bookkeeping: name the field, mark membership, order
consumption, and state the residual only where the field authorizes the read.

Support-local quiet also prevents a common false picture of computation. The
apparatus does not begin with an infinite mathematical object and then trim it
down for convenience. It begins with a finite record that can circulate
through compatible events. The record contains candidates, not a vague cloud
of possible candidates. The trace orders those candidates. The residual tests
those candidates. The result closes only where the record let the apparatus
read. This is not a weakness in the theory. It is the price of a witness that
can actually be held.

The multiset language helps because the field may care about repeated
presence without yet caring about order. Two copies of a candidate can matter
as two carried chances, two receipts, or two positions in a later accounting,
even before the trace decides which copy comes first. A plain set would erase
that multiplicity too soon. A sequence would impose order too soon. The finite
field therefore keeps exactly the information this section needs: which
candidates are in scope, with whatever finite multiplicity the apparatus has
chosen to carry, before any execution order begins.

Once execution begins, the trace converts the unordered field into a history.
At each step it chooses a supported candidate, spends that choice, and carries
the remainder. The history is not merely a list of names. It records the
finite cost of selection, the shrinking field of alternatives, and the place
where residual quiet can later attach. By the end of the trace, the certificate
can say: this was the support; this was the order in which the support was
consumed; this is the residual statement earned inside that support.

This order of explanation also protects the next section. Partitioned real
coordinates will soon give the apparatus a way to name supported candidates in
a coordinate frame. That step should not retroactively turn the present
support into an interval, a cell decomposition of a continuum, or an analytic
estimate. Coordinates will help name the finite field. They will not supply a
secret infinity behind it. The present section must therefore stop at the
support/order grammar: membership before execution, consumption with decrease,
and residual quiet for supported variations.

The doctrine can be stated compactly. A finite support \(K\) names the
candidates this certificate may read. An execution order consumes candidates
from \(K\) one finite event at a time. Each consumption event selects a member,
keeps the remainder inside the original field, and decreases the field. A
closed trace may report \(r_{\mathrm{term}}(b')=0\), and where the spline
identification applies it may report \(r_{\mathrm{spline}}(b')=0\), only for
\(b'\) in \(K\). The apparatus has not read every possible middle. It has made
one finite field tangible, carried it through an order, and closed the
residual on the candidates that field allowed it to read.

\section{Coordinates Without Imported Analysis}

Once a support has named the candidates in scope, the apparatus still needs a
way to name their positions. The trace from the previous section can say which
candidate it consumes next, and it can keep residual quiet local to the
carried field. But a later calculation will also need coordinate-facing
speech: before, after, adjacent, bounded step, and located candidate. This
does not require the book to borrow real analysis. It requires a finite
coordinate contract.

The contract begins with an ordered axis. Publicly, write its named points as
\[
  x_0 \le x_1 \le \cdots \le x_N .
\]
The display says only that the apparatus has finitely many coordinate names
and an order among adjacent names. It does not claim that the axis is a
completed continuum. It does not assert a normed function space or an
analytic theorem about limits. The coordinate system gives the finite support
a language of position. That is already a substantial act: the apparatus can
now say where a supported candidate stands relative to the other names it
carries.

A partition adds a finite scale to this order. Between adjacent names the
apparatus records a bound
\[
  d(x_i,x_{i+1}) \le h .
\]
Here \(d\) is the apparatus distance and \(h\) is the mesh bound. The bound
does not promise convergence. It does not say that shrinking \(h\) recovers a
continuum object. It says that adjacent coordinate names in this finite
partition stand within a controlled step. The mesh therefore measures a local
adjacency in the carried partition, not a global analytic limit.

The finite field also needs an admissible range. The apparatus does not assign
coordinates to every possible body or every possible candidate. It begins with
a carried body and a finite candidate support, and it asks whether those
items lie inside the allowed range. The body must be admitted. Each supported
candidate must be admitted. Only then does the partition become a coordinate
record for the present data. This keeps the support from pretending to be a
universal domain: it remains a finite field that has passed a range check.

Location is the next act. A supported candidate receives a coordinate name:
\[
  b \in K
  \quad \Longrightarrow \quad
  x(b) \in \{x_0,\ldots,x_N\}.
\]
The implication matters. The candidate receives a coordinate name because it
belongs to \(K\). The partition does not coordinate unsupported candidates by
silence, and it does not discover candidates outside the field. It makes the
carried support tangible in a coordinate grammar. The record can now say:
this candidate belongs to the support, this candidate lies inside the
admissible range, and this candidate has this coordinate name among the
finite names of the partition.

The trace must keep faith with that coordinate record. Its knots cannot drift
away from the carried data. If the support from the trace and the support from
the partitioned data are being identified, the public summary may be written
as
\[
  K_{\mathrm{trace}} = K_{\mathrm{data}} .
\]
This is not a new theorem of the continuum. It is bookkeeping for the finite
certificate. The knots the trace consumes are the candidates the partition
locates. The order still belongs to the trace; the coordinate names belong to
the partition; the support keeps them talking about the same finite field.

Residual quiet transfers through that match. When a terminal residual is zero
on the trace knots, the partitioned data may read the same quiet on its
supported candidates. If weak language is useful, the statement should remain
local:
\[
  r_{\mathrm{weak}}(b')=0
  \qquad \text{for } b' \in K .
\]
And, where the spline identification from the earlier stationarity chapter is
being carried, the same local discipline gives
\[
  r_{\mathrm{spline}}(b')=0
  \qquad \text{for } b' \in K .
\]
The phrase "for \(b' \in K\)" does the work. It prevents the coordinate
partition from inflating the claim into global variation control. The weak
residual vanishes for supported candidates. The spline residual vanishes for
supported candidates. Unsupported candidates remain outside this certificate's
voice.

The word "weak" also needs a boundary. In this section it names a
partition-facing pairing language for finite functions generated by the
support. It does not announce completed Sobolev analysis. It does not supply an
ambient function-space theorem, inner-product theorem, or compactness
argument.
The apparatus has a finite partition, a finite pairing discipline, and a
support-generated family of functions. That is enough for the present
coordinate-facing claim. It is not enough to claim a completed theory of weak
solutions.

This restraint keeps the chapter moving in the right order. The first section
made support prior to execution. This section gives that support coordinate
names, adjacency, a mesh bound, admissibility, and knot matching. The next
section may use those ingredients to write the next algebraic laws.
It should not inherit an unstated analytic universe from this one. Coordinates
prepare the finite field for algebra; they do not solve the algebra in
advance.

The result is a sharper finite promise. A support says which candidates the
certificate may read. A coordinate partition says where those supported
candidates stand among finitely many ordered names and how adjacent names are
bounded by the mesh. A knot match says that the trace consumes the same field
the partition locates. Residual quiet then remains exactly where it belongs:
on the supported candidates, in the coordinate-facing language the apparatus
has actually carried.

\section{Cubic Middle-Node Equations}

The coordinate contract prepares the finite field for equations. It names
supported candidates, places them among finitely many ordered coordinates, and
keeps residual quiet local to the support that the apparatus actually carries.
The next question is no longer whether a candidate has a name. It is what
condition a named interior candidate must satisfy once the apparatus asks a
path to be quiet. The answer begins with the first case where there are two
interior names at once.

Consider a cubic support written as a four-node path
\[
  q_0 \longrightarrow q_1 \longrightarrow q_2 \longrightarrow q_3 .
\]
The endpoints \(q_0\) and \(q_3\) are held fixed for the present certificate.
They are the boundary of the local calculation. The middle names \(q_1\) and
\(q_2\) are the supported interior records. The word "middle" does not mean
that the book has imported a continuum point between two completed endpoints.
It means that the apparatus has carried two named interior positions in the
finite record. Each middle name stands only because a support has made it
available and a coordinate contract has made its location tangible.

This is the first reason the cubic case is useful. A single middle node can
teach a local balance law, but it can hide the fact that a path with two
interior names has more than one local question. The left middle name may be
tested while the right middle name is held fixed. The right middle name may be
tested while the left middle name is held fixed. And the two middle names may
move together, in which case the segment between them changes from both sides.
The section must keep these three readings separate. Otherwise a finite
calculation would quietly turn two supported local equations into an
unsupported claim of whole-path quiet.

The finite record also decides what kind of object the equations are allowed
to be. The four names are not samples from a curve that has already been
granted. They are the names the certificate has in hand. The two boundary
names let the local problem stand still long enough to be measured. The two
middle names are the places where the apparatus may ask for balance. The
equations therefore begin as accounting over named positions. They do not
begin as a law over an unbounded space of paths and then shrink down to the
finite case. Measurement runs in the other direction: name the finite field
first, then state the condition that the field can actually test.

The action for the four-node path is the sum of three segment costs:
\[
  S(q_0,q_1,q_2,q_3)
    = C(q_0,q_1) + C(q_1,q_2) + C(q_2,q_3).
\]
The first term reads the segment from the fixed initial endpoint to the left
middle name. The second term reads the segment between the two middle names. The
third term reads the segment from the right middle name to the fixed target.
This display is deliberately finite. It does not ask for a differentiable
curve, a completed function space, or a background continuum of variations. It
asks the apparatus to add the three costs that its record has made available.

The equality also records what has not yet been decided. The cost of the whole
four-node path is not a single opaque number handed down from outside the
apparatus. It is assembled from the three finite reads the record permits.
That assembly matters because each middle name touches two segments. The left
middle name touches the first and second segment. The right middle name
touches the second and third segment. Thus the middle segment is shared. It is
the place where the two interior names can no longer be treated as completely
independent once they move together.

The notation \(S\) should therefore be read as a carried sum, not as a hidden
functional over all possible paths. It is large enough to ask whether the
supported middle names are quiet relative to the three segment costs. It is
not large enough to speak for unsupported middle names or for a path that has
not been named by this certificate. The support still governs the voice of the
equation. The coordinate contract still governs where the named candidates
stand. The cost only measures the finite path once those acts have already
been performed.

Begin with the left middle name. Hold \(q_0\), \(q_2\), and \(q_3\) fixed, and
let the apparatus test a supported left variation \(\eta_L\) of \(q_1\). The
left test changes the first segment and the middle segment, because \(q_1\)
belongs to both. It does not change the right endpoint segment from \(q_2\) to
\(q_3\). The finite residual recorded by this one-coordinate test is written
as \(R_L(\eta_L)\). The left middle-node equation is
\[
  R_L(\eta_L)=0
  \qquad \text{for supported left variations }\eta_L .
\]
The last phrase is part of the equation. It is not an explanatory decoration.
The residual vanishes on the supported left variations that the certificate
has admitted. If a left variation is not in the carried support, this
certificate has not tested it.

This left equation is already an Euler-Lagrange-shaped statement, but the
shape must be read with finite discipline. The apparatus has not taken a
derivative in a completed space of curves. It has not ranged over arbitrary
perturbations of a continuum path. It has compared a named middle record with
the supported alternatives the trace can read. The finite Euler-Lagrange
equation says that the left middle name is locally balanced against the
supported left tests. It is a residual equation on the named support.

The word "residual" keeps the claim from becoming merely verbal. A residual is
what remains when the apparatus compares the present middle name with the
supported alternatives available to it. If the residual is nonzero, the
support has found a local imbalance. If the residual is zero for every
supported left test, the apparatus has no left-supported discrepancy left to
spend in this one-coordinate reading. That is a real closure, but it is a
closure with an address: the left support, the left middle name, and the
fixed surrounding nodes of the present cubic record.

The same discipline applies on the right. Hold \(q_0\), \(q_1\), and \(q_3\)
fixed, and test a supported right variation \(\eta_R\) of \(q_2\). This
variation changes the middle segment and the final segment. It does not
change the first segment from \(q_0\) to \(q_1\). The right residual is written
\(R_R(\eta_R)\), and the right middle-node equation is
\[
  R_R(\eta_R)=0
  \qquad \text{for supported right variations }\eta_R .
\]
Again the support condition is not optional. The equation does not say that
the path is quiet against every rightward motion that a later theory might
invent. It says that, inside the present finite record, every admitted right
test has zero residual.

The right equation also prevents a false symmetry. Since both middle names
touch the central segment, one might be tempted to treat the two balances as
one undifferentiated interior condition. The apparatus refuses that shortcut.
The left test asks what happens when the left interior name is replaced by a
supported left alternative. The right test asks what happens when the right
interior name is replaced by a supported right alternative. They share part of
the cost, but they are different tests with different supports. The finite
record has to keep both receipts.

The two equations are parallel but not interchangeable. The left equation
looks at the cost terms touched by \(q_1\). The right equation looks at the
cost terms touched by \(q_2\). Both see the middle segment, but they see it
from different one-coordinate tests. This is why the cubic path is not merely
a longer version of the one-middle case. It forces the apparatus to carry two
interior balances without collapsing them into one anonymous interior
condition.

The phrase "finite middle-node balance" is useful here. A balance is a
relation among costs the apparatus can actually compare. The left balance
compares the supported alternatives for the left middle name while the rest of
the cubic record remains fixed. The right balance compares the supported
alternatives for the right middle name under the matching fixed-boundary
condition. The apparatus earns two balances because it has two supported
middle names. It does not earn a global variational principle by saying the
word "balance" twice.

The finite Euler-Lagrange boundary is therefore exact but narrow. It is exact
because the residual statements are not metaphors. They are the algebraic
conditions the carried cubic record must satisfy at its supported middle
names. It is narrow because every quantifier remains inside the admitted
support. The book may call these finite Euler-Lagrange equations, provided the
the adjective must stay strict. Finite means the equation is tested on named
variations, inside the field the apparatus has carried, with endpoints and the
other middle name held according to the one-coordinate test.

This boundary is not a hedge. It is the content of the construction. A
completed analytic setting may have its own way to discuss variation, but this
chapter is not leaning on that setting. It is showing what the apparatus can
say before such a setting is available. The finite Euler-Lagrange equation is
the statement that no supported one-coordinate test exposes residual pressure
at the middle name being tested. The equation is strong because it is exact on
the support. It is limited because the support is the place where the
apparatus has earned a voice.

This language also preserves the role of the endpoints. Fixed endpoints do
not become metaphysical absolutes. They are fixed for this local calculation.
The certificate is asking what the two supported interior names must do while
the boundary names stand still. A later certificate may move a boundary or
compare a family of boundaries, but that is a new question. Here the endpoints
define the local field of the cubic test. They make the middle equations
legible by preventing the whole path from sliding at once.

The same point can be put as a bookkeeping rule. A middle-node equation is not
a declaration about the identity of the whole path. It is a statement about
what changes when exactly one supported interior name is exchanged while the
rest of the local record is held. This is why the equations can be written
before the chapter has introduced a selector. The selector will later ask
which candidate survives a measured choice. The present section only supplies
the local tests that candidates must face.

Now let both middle names move together. The apparatus tests a pair
\((\eta_L,\eta_R)\) and compares the cost of the changed cubic path with the
cost of the original one. The resulting coupled difference is not merely the
sum of the two one-coordinate residuals. The middle segment has been touched
from both sides at once. When \(q_1\) and \(q_2\) both change, the cost
\(C(q_1,q_2)\) is replaced by the cost of the changed middle segment, and that
replacement contains a contribution that neither one-coordinate test sees by
itself.

Write the coupled accounting as
\[
  \Delta_{12}S(\eta_L,\eta_R)
    = R_L(\eta_L) + R_R(\eta_R) + M(\eta_L,\eta_R).
\]
The term \(M(\eta_L,\eta_R)\) is the mixed coupling. It records the extra
middle-segment contribution that appears only when the two supported middle
names move together. The display is the chief warning of this section. The
left residual and the right residual can both vanish on their respective
supports, and the coupled variation can still contain a mixed term. The
apparatus must not hide that term inside the two residual equations.

The mixed term is not an error in the finite account. It is the honest residue
of asking a two-middle question. One-coordinate variation isolates a middle
name by holding the other one fixed. Coupled variation removes that protection.
The middle segment then depends on both changed names, and the accounting must
say so. The finite theory becomes stronger, not weaker, when it refuses to
erase this residue. It shows exactly where the local equations stop and where
an additional coupling condition begins.

This is also where the language of two equations becomes most precise. There
are two one-coordinate residual equations, not because the apparatus has split
a single smooth condition into pieces, but because the record contains two
supported interior names. Their coupled movement is an additional finite
question. If the mixed term is present, the apparatus has learned something:
the two local quiet statements have not exhausted the relation between the
middle names. The middle segment still carries pressure that belongs to the
pair.

The clean exact closure is
\[
  R_L=0,\quad R_R=0,\quad M=0
    \quad\Longrightarrow\quad
    \Delta_{12}S=0.
\]
This implication should be read with the same support boundary as before. The
left residual is zero on supported left variations. The right residual is zero
on supported right variations. The mixed term is separately solved, cancelled,
or otherwise controlled for the coupled supported test. Only then does the
coupled difference vanish. The two middle equations alone do not erase the
coupled variation.

The implication also separates exact closure from later control. In the clean
display, the mixed term is zero. In prose, the chapter may leave room for a
later finite argument that bounds or carries the mixed contribution as
residue. Those are not the same claim. Zero says the coupled contribution has
closed in the present calculation. A bound says a later accounting may keep
the residue under control. A carried residue says the apparatus knows the
pressure remains. Here the exact case teaches the algebra; broader control
belongs to the later sections on selection and squeezing.

This is the exact place where the section must resist a tempting shortcut. It
would be easy to say that the left and right equations make the cubic path
stationary. That sentence is too large. The equations make the supported
one-coordinate residuals quiet. They do not, by themselves, prove quiet for
all coupled movement. They do not range over unsupported variations. They do
not smuggle in a smoothness argument that declares the mixed term negligible.
They do not import a completed variational calculus. The coupled statement
needs the mixed term to be named and handled.

The terminology of residue is appropriate here. After the two middle-node
equations have done their work, the mixed coupling remains as the next item of
accounting. It may be zero in a special exact case. It may be cancelled by a
separate condition. It may be bounded by a later finite argument. But it must
not disappear because the prose wants the cubic path to look simpler. A
measurement record is reliable only when the remaining pressure is still
visible.

The mixed term also protects the meaning of "quiet." Quiet at the left middle
name means no supported left test reports residual pressure. Quiet at the
right middle name means no supported right test reports residual pressure.
Coupled quiet is a stronger sentence. It requires the paired movement to leave
no additional middle-segment pressure beyond the two residuals. The apparatus
therefore has a hierarchy of claims: left quiet, right quiet, and coupled
quiet. Confusing the hierarchy would make the record sound more complete than
it is.

This also clarifies what the cubic equations contribute to selection. They do
not yet choose one candidate from a support. They define the local equations
that a supported interior pair must satisfy before selection can be a
measured act rather than a preference. A candidate pair that fails the left
residual test has not met the left middle-node balance. A candidate pair that
fails the right residual test has not met the right middle-node balance. A
candidate pair that passes both tests may still carry coupled residue through
the mixed term. The apparatus now has a sharper set of questions to ask of
the support.

This sharpening is the section's practical contribution to the chapter. The
previous section made finite coordinates available without importing analysis.
This section turns those coordinates into equations that selection can read.
The equations are not yet a choice procedure. They are tests. A later
selection rule can prefer, reject, or compare supported candidates only after
the residual tests have given the candidates a measured profile. The cubic
case gives that profile two middle-node faces and one coupled residue.

The word "local" should therefore remain visible. The left equation is local
to the left middle support. The right equation is local to the right middle
support. The mixed term is local to the coupled movement of the two supported
middle names. None of these statements surveys an unbounded landscape of
paths. They organize the finite record already in hand. They tell the
apparatus where quiet has been earned and where an unresolved coupling still
stands.

The result can be summarized in one sentence. A cubic support with fixed
endpoints and two supported middle names carries two finite Euler-Lagrange
equations, one for each middle name, and a separate mixed coupling residue
when the two middle names move together. The section has not completed the
selection problem. It has prepared it. The next step can ask how an apparatus
chooses among locally measured candidates without confusing selection with
global uniqueness. A later step can ask how squeezed supports carry residual
behavior forward. For now the boundary is exact: two middle-node residuals,
one visible mixed term, and no claim beyond the support the certificate has
actually named.

\section{Heartbeat Selection as Local Measured Choice}

The cubic equations give the apparatus local tests. They say when a supported
left middle name is quiet, when a supported right middle name is quiet, and
where a mixed coupling remains if the two middle names move together. They do
not yet say which admitted candidate the apparatus reads next. A support may
contain several candidates that satisfy the current local tests, or several
candidates that still need to be tested. Selection is the next discipline: it
turns a finite field of admitted candidates into an ordered act of reading.

The support remains prior. Selection does not create a candidate, discover an
outside option, or enlarge the multiset that the certificate carries. It works
only on candidates already admitted. The new ingredient is a heartbeat cost:
a measured priority attached to each admitted candidate. Publicly, write
\[
  h:K\to \mathbb{N}.
\]
Here \(K\) is the current finite support and \(h(b)\) is the heartbeat cost of
the admitted candidate \(b\). The codomain \(\mathbb{N}\) matters. The
heartbeat is a finite number the apparatus can compare. It is not a biological
pulse, not a runtime measurement, and not a probability assigned to a possible
world. It is an ordering cost carried by the present record.

This keeps the word "heartbeat" from becoming decorative. The heartbeat is a
clock only in the finite sense that it gives the apparatus a countable priority
scale. Two candidates can be compared because their costs are numbers. A
candidate with cost \(2\) stands before a candidate with cost \(5\) under this
policy. The apparatus has not measured duration in a machine, and it has not
found a secret physical pulse inside the candidate. It has attached a finite
priority to an admitted name so that the next read can be justified by a
carried comparison.

This cost lets the trace ask for a fastest admitted read. A candidate
\(b_\ast\) is heartbeat-minimal when
\[
  b_\ast \in K,
  \qquad
  h(b_\ast)\le h(b)
  \quad\text{for every } b\in K .
\]
The first clause says that the candidate belongs to the current support. The
second says that no admitted candidate has a smaller heartbeat cost. This is
minimality, not uniqueness. If two admitted candidates have the same least
cost, both may be heartbeat-minimal. The finite support still has a measured
front edge, but it has not necessarily collapsed to a single candidate by cost
alone.

That boundary is important enough to keep in the prose. A nonempty finite
support has at least one heartbeat-minimal admitted candidate. This gives the
apparatus existence of a next measured read. It does not give the stronger
claim that there is exactly one minimizer. If the apparatus needs one actual
next candidate and several candidates tie, it must carry a tie policy, or it
must choose one minimizer as an additional witnessed act. The heartbeat
orders candidates by cost. It does not erase every ambiguity that finite cost
can leave behind.

The tie case is not a defect. It is a visible boundary in the record. Suppose
two supported candidates both have the least heartbeat cost. The cost policy
has done its work: it has ruled out every candidate with larger cost as the
next fastest read. What remains is a finite front of tied minimizers. To move
from that front to one consumed candidate, the apparatus needs a secondary
rule, a stable convention, or a witnessed choice. That extra act should be
named when it is used. Otherwise the prose would pretend that a preorder had
become a single-valued selector by silence.

The measured choice step is therefore written as
\[
  K \xrightarrow{\,b_\ast\,}_{h} K',
  \qquad
  b_\ast\in K,
  \qquad
  K'\subseteq K,
  \qquad
  |K'|<|K|.
\]
Read the display as fastest admitted consumption. The apparatus begins with
the finite support \(K\). It chooses a heartbeat-minimal admitted candidate
\(b_\ast\). It consumes that candidate and carries a remaining support \(K'\).
The remainder stays inside the original support, and the support strictly
decreases. The heartbeat subscript records the policy by which the next read
was chosen; it does not change the membership rule.

The decrease clause keeps the trace honest. If \(K'\) were allowed to have the
same size as \(K\), the apparatus could announce a fastest read without
spending any support. If \(K'\) were allowed to contain candidates outside
\(K\), the heartbeat would become a hidden support generator. The display
permits neither. The consumed candidate comes from the present support, and
the remainder stays inside the present support while becoming strictly
smaller. Selection is therefore an event of finite expenditure, not a fresh
act of creation.

This is the local measured choice promised by the title. It is local because
the comparison is made only inside the current support. It is measured because
the comparison uses the heartbeat cost. It is choice because the apparatus
commits to a next read, including any tie policy or witnessed choice needed
when several candidates are equally minimal. It is not global optimization.
The apparatus has not searched all possible candidates. It has not minimized
over a continuum, over a hidden background space, or over candidates that the
support did not admit.

The heartbeat also does not replace the execution trace. It refines the order
in which the trace consumes the support. Once the heartbeat-ordered act has
chosen and consumed its candidate, the event can still be read as an ordinary
ordered trace: one admitted candidate was consumed, and the remaining support
was carried forward. Forgetting the heartbeat order leaves the support/order
discipline from the opening section intact. The heartbeat tells why this
candidate was read now; the ordinary trace still records that a supported
candidate was read and that the support became smaller.

This forgetting is not a loss of truth. It is a change of view. With the
heartbeat visible, the record says that the candidate was read because it was
minimal under \(h\). With the heartbeat forgotten, the record still says that
the candidate belonged to the support and that the trace consumed it. The
ordinary trace is enough for residual bookkeeping that only needs membership
and consumption. The heartbeat order adds a measurable reason for the
itinerary, then allows the older support grammar to keep doing its work.

This forgetting is useful because residual quiet was already formulated for
the ordinary trace. The heartbeat-ordered trace does not need to invent a new
kind of residual. It carries the same terminal and support-local residual
statements through a more disciplined order of consumption:
\[
  r_{\mathrm{term}}(b)=0
  \qquad\text{for } b\in K .
\]
The phrase "for \(b\in K\)" remains the guardrail. Heartbeat order has not
expanded the domain of the residual. It has only supplied a measured policy
for deciding which admitted candidate is consumed next.

The partitioned-coordinate story also carries forward. A heartbeat-ordered
trace can still be read by the partitioned approximation once the support and
the coordinate record match. The finite functions and residual statements do
not become analytic claims merely because the trace now has a cost policy.
They remain support-local. The heartbeat supplies a measurable order on the
admitted candidates; the partition supplies coordinate names; the residual
reports quiet on the supported tests the certificate has earned.

Thus the section has three layers and should not collapse them. Support says
which candidates can be read. Coordinates say where those supported candidates
stand. Heartbeat selection says which admitted candidate is read next under a
finite priority. None of the three layers replaces the others. If the support
is missing, the heartbeat has no lawful domain. If the coordinate record is
missing, the candidate may be admitted but not yet located in the finite
coordinate grammar. If the heartbeat is missing, the trace may still consume
the support, but it lacks this measured priority policy.

The cubic equations from the previous section should enter here only as
carried structure. The section may note that, on supported tests,
\[
  R_L(\eta_L)=0,\qquad R_R(\eta_R)=0
  \qquad\text{on supported tests.}
\]
But it should not redo the cubic calculation. The heartbeat does not solve the
mixed coupling. It does not prove coupled quiet. It does not change the two
middle-node equations into a global theorem. It says how a finite apparatus,
already holding local equations and residual boundaries, chooses which
admitted candidate to read next.

This distinction prevents a second overclaim. Selection is not the same thing
as a uniqueness guarantee. A fastest admitted read may be unique after a tie
policy, or after the support happens to contain a single minimal candidate, or
after the apparatus witnesses one minimizer among several. Those are different
facts. The heartbeat cost alone supplies a preorder by finite number. It gives
the apparatus a measured front of the support. The actual next read is the
minimal candidate the apparatus carries forward.

Choice should therefore be heard in its measured sense. The apparatus
does not choose because it prefers a candidate in silence. It chooses because
the support is finite, the heartbeat cost is carried, and a minimal admitted
candidate has been selected for consumption. If there are ties, the tie policy
or chosen witness is part of the measurement record. If there are no ties, the
heartbeat cost itself has isolated the next read. In both cases the claim is
local to the current support.

Nor is heartbeat selection a probabilistic rule. Nothing in the display says
that \(h(b)\) is a likelihood, a frequency, or a weight of belief. The number
is a cost used for ordering consumption. A candidate with a lower heartbeat is
read earlier under this policy; it is not thereby more true, more probable, or
more real than another supported candidate. Truth and residual quiet remain
attached to the certificate's residual statements. The heartbeat only orders
the finite work by which the certificate spends its support.

The remaining support \(K'\) keeps failure and continuation honest. If
unread candidates remain, the apparatus has not forgotten them. It carries
them forward as the finite field for the next measured choice. If the support
is exhausted, the trace has spent the candidates it was authorized to read.
At no point does the heartbeat permit the trace to replenish the support from
outside. A new candidate requires a new support act. A new support act is not
hidden inside the word "fastest."

The local doctrine can therefore be stated compactly. A finite support \(K\)
names the admitted candidates. A heartbeat cost \(h\) assigns a natural-number
priority to those candidates. A heartbeat-minimal candidate has no admitted
candidate below it in cost. The trace consumes one such candidate, records the
remaining support, and can forget the heartbeat policy back to the ordinary
ordered trace. Residual quiet remains support-local, and the cubic middle-node
equations remain carried tests rather than a new selection principle.

This doctrine also explains why the section does not need an external
optimization principle. The apparatus is not borrowing a general theory of
best choice. It is performing a finite comparison among names already in hand.
The whole grammar remains inspectable: a support \(K\), a cost assignment
\(h\), a minimal admitted candidate \(b_\ast\), a consumption event, and a
remaining support \(K'\). Nothing else is hidden behind the notation.

This is where the section should stop. It has given the apparatus a measured
choice policy for the current finite support. It has explained ties without
pretending they vanish. It has carried residual quiet through heartbeat order
without hiding the mixed coupling from the cubic calculation. The next
question is different: how a support can be squeezed while residual behavior
is preserved or controlled. That is not a selection rule. It is a later
question about families of supports and the residue they carry.

\section{Squeezed Residuals in a Finite Tower}

Heartbeat selection spends one finite support in a measured order. It chooses
inside the support already admitted, consumes one candidate, and carries the
remaining finite support forward. That doctrine leaves a second question for
the close of the chapter. A single support can be read, but the apparatus may
also compare a whole family of supports, each still finite, each refining the
last, each carrying a residual record. The issue is no longer which candidate
comes next inside one support. The issue is how residual size behaves as the
apparatus passes through a tower of finite supports.

Write the stages as \(K_0,K_1,K_2,\ldots\). The notation should not suggest an
infinite support at any one stage. Each \(K_n\) is finite. Each stage is an
admitted record the apparatus can hold, inspect, and pass forward. Refinement
means that what a stage has already admitted does not fall out silently at the
next stage. A carried candidate remains carried, or else the record must name
the loss as a new boundary. In the ordinary case considered here, the tower
keeps prior admitted contents available while it adds resolution.

This is the tower version of finite honesty. The apparatus does not announce a
continuum merely because it can name many finite stages. It names a sequence
of finite stages and a rule by which the stages relate. The rule matters as
much as the stages. Without a refinement rule, the family would be a list of
unrelated supports. With the rule, the family becomes a tower: a finite record
at each level, and a visible passage from one level to the next.

The passage from \(K_n\) to \(K_{n+1}\) should be read as a receipt-preserving
event. What \(K_n\) made tangible for the apparatus remains available to
\(K_{n+1}\), unless the next record explicitly marks a boundary where
availability stops. This is how refinement differs from replacement. A
replacement can abandon one support and begin again. A refinement carries the
earlier support into a better-resolved stage. The tower therefore lets the
apparatus compare residual reports across stages without smuggling in a hidden
change of subject.

That receipt-preserving structure also explains why the tower may speak about
approximants. An approximant at stage \(n\) is not a ghost from an outside
space. It is a name accepted by the stage support. A sequence of approximants
is therefore a sequence of finite receipts, one per stage, not a single
unbounded object disguised as notation. The tower makes those receipts
fungible across refinement: each stage can pass its admitted information to
the next stage while still remaining a finite stage.

The word "completion" enters only under that discipline. A completion claim
does not say that Chapter 7 has proved a completed analytic space. It says
that the tower comes with a declared relation by which an approximating
sequence may be accepted as converging to a limit name. The limit is not
outside the declared completion, because the completion relation itself says
how the approximants carried by the tower point to it. That is a membership
sentence inside the declared relation, not an unqualified continuum theorem.

This distinction protects the chapter boundary. The tower may prepare a door,
but it does not walk through every room beyond it. Chapter 7 may say that a
limit name belongs by declaration to the completion carried by the tower. It
may not say that all finite geometry has become a completed analytic geometry,
or that the later variational theory has already been established. Those are
later questions. Here the completion relation serves one purpose: it prevents
the limit from becoming a mysterious outside object while also
preventing the prose from claiming more analysis than the finite apparatus has
earned.

Residuals now need a measured size language. A terminal residual may begin as
an integer or a finite report attached to a stage trace, but convergence talk
requires a magnitude. Let \(r_n(\eta)\) name the magnitude assigned at stage
\(n\) to the residual tested by \(\eta\). The symbol \(\eta\) still denotes an
admitted test; it does not range over every possible analytic perturbation.
The magnitude language lets the apparatus compare residual size at different
stages and ask whether the sizes become arbitrarily small under the admitted
tolerance relation.

The key assertion is the squeeze:
\[
  0 \le r_n(\eta) \le B_n(\eta),
  \qquad
  B_n(\eta)\to 0 .
\]
The lower inequality says that residual magnitude is measured as a nonnegative
quantity. The upper inequality says that the residual at stage \(n\) never
exceeds a bound \(B_n(\eta)\). The last clause says that this bound tends to
zero along the tower. The apparatus therefore need not compute a separate
limit for the residual. It can compare the residual to a vanishing bound and
let the bound carry the limiting behavior.

The consequence is the expected one:
\[
  r_n(\eta)\to 0 .
\]
This statement remains pointwise in the admitted test \(\eta\). It does not
assert uniform control over every possible test, and it does not transform the
finite tower into a completed continuum. It says that, for the test under
discussion, the residual magnitude is squeezed below a bound that vanishes.
The claim is strong enough to close the finite-support arc and modest enough
to leave later analytic structure untouched.

The bound may come from a density certificate. In ordinary mathematical
language one might invoke a Weierstrass-style promise: a sufficiently rich
family can approximate the target with a bound that shrinks. In this book that
promise should be read as an admitted certificate, not as a theorem rederived
inside the chapter. The certificate supplies \(B_n(\eta)\). Once the tower
carries both the residual and the bound, the squeeze does the local work:
residual below bound, bound to zero, residual to zero.

The certificate has a precise burden. It does not need to tell the whole story
of approximation. It needs to provide, for the admitted test, a bound that the
residual magnitude can sit below and that the tower can make smaller beyond
any accepted tolerance. Once that burden is met, the apparatus has enough to
carry a finite limiting claim about residual magnitude. It does not need to
identify a preferred analytic model, and it does not need to describe every
possible way of refining the support. The finite tower and the bound do the
work that this chapter needs.

This also keeps the squeeze from becoming a metaphor without count. The
display has three inspectable parts: nonnegative residual magnitude, residual
below the bound, and a bound that vanishes along the tower. Remove any one of
the three and the claim changes. Without nonnegativity, "size" loses its
order. Without the upper bound, the residual no longer rides the certificate.
Without the vanishing promise, the bound may control the residual while never
forcing it down. The displayed inequality is therefore the grammar of the
section, not an ornament.

Two witness families can carry the same shape of argument. A B-spline witness
and an FFT witness need not be the same basis. They do not erase basis choice,
and they do not prove an all-bases theorem. They show that different admitted
witness types may present the same residual grammar: a finite tower, a
residual magnitude, a bound, and a vanishing promise for the bound. The public
point is the shared shape of the certificate, not identity of the bases.

This is why the section should avoid a stronger headline claim about basis
neutrality. The apparatus has not surveyed every possible basis and found the
same theorem waiting there. It has received witness types that can be
translated into the same squeeze form. The witness matters. Basis choice still
belongs to the record. What circulates across the witnesses is the grammar of
the squeeze: if the residual lies below an admitted vanishing bound, the
residual magnitude tends to zero along the finite tower.

The stage trace connects this magnitude statement back to the terminal
residuals from the earlier sections. The tower does not invent a second
residual unrelated to the trace. It measures the terminal residual at stage
\(n\):
\[
  r_n(\eta)
  = \mu\!\left(r_{\mathrm{term},n}(\eta)\right),
  \qquad
  r_n(\eta)\to 0 .
\]
Here \(\mu\) denotes the magnitude reading of the terminal residual. The
display says that the residual magnitude used in the squeeze is the measured
size of the stage terminal residual. The terminal residual remains the same
kind of finite report as before; the magnitude reading makes it comparable
along the tower.

The magnitude reading also prevents a category mistake. The terminal residual
itself belongs to the finite trace: it is the report left at the end of a
stage computation. The squeeze needs a magnitude because it compares sizes
across stages and against a tolerance. The symbol \(\mu\) marks that passage
from a terminal report to a size the tower can compare. It does not add a new
residual, and it does not erase the finite origin of the report. It makes the
terminal report tangible as a measured magnitude the tower can carry.

Carried knots still have their sharper local statement:
\[
  \eta\in K_n
  \quad\Longrightarrow\quad
  r_{\mathrm{term},n}(\eta)=0 .
\]
If the test belongs to the stage support, the stage trace solves that
terminal residual outright. This is stronger than smallness at that stage,
but it is local to carried support. The squeeze statement handles residual
size along the tower. The carried-knot statement handles exact terminal quiet
where membership at the current stage has been earned. The two claims should
not be merged.

This separation keeps the chapter's residual doctrine clean. Membership in a
stage support gives exact local quiet for the terminal residual named by that
membership. A vanishing bound gives limiting smallness for residual magnitude
along the tower. Declared completion says that the limit name is not outside
the completion relation carried by the tower. These are three distinct
sentences: local exactness, squeezed smallness, and declared completion
membership. The book needs all three, but it should not let one pretend to be
the others.

The finite nature of the stages also remains visible after the limit language
appears. At no point does the tower become a single infinite support that the
apparatus reads all at once. It remains a sequence of finite supports. The
limit statement speaks about behavior along the sequence. The completion
relation records how the approximants point to a limit name. The support at
stage \(n\) is still \(K_n\), and the residual statement at stage \(n\) still
belongs to the tests admitted at that stage or measured against that stage's
bound.

The apparatus should therefore keep two temporal directions apart. Within one
stage, the apparatus asks whether a supported test has terminal quiet. Along
the tower, the apparatus asks whether residual magnitudes shrink under the
admitted bound. The first direction is local and exact. The second direction
is staged and limiting. The section needs both because a carried knot can be
quiet at its stage while other admitted tests are controlled only by the
vanishing bound. Refinement does not flatten those two questions into one.

This is the final movement of Chapter 7's numbered arc. The chapter began by
requiring finite support before measurement could say anything definite. It
then counted the support, gave it finite coordinates, extracted local cubic
equations, and ordered a current support by heartbeat choice. The present
section adds one more operation: arrange finite supports into a refining
tower and squeeze residual magnitude by an explicit vanishing bound. The
apparatus can now speak about residual behavior along refinement without
pretending that a single finite stage already contains the continuum.

The boundary is equally important. Nothing in this section introduces an
analytic completion theorem for the later variational theory. Nothing here
says that all bases have disappeared, that all tests are controlled uniformly,
or that the finite tower has become the full geometry of a completed space.
The chapter closes with residual towers, not with the later theory of forms
and completion. It has prepared a completion-facing door by tanging the finite
records and funging their residual reports through an explicit squeeze, but it
has not yet crossed into the later theory that studies what stands beyond that
door.

\section*{Bridge: Residual Towers Ask for the Invariant}

Chapter 7 has earned a controlled remainder, not a governing measure of that
remainder. Support gave measurement a finite field; coordinates tanged that
field into named positions; local cubic equations made residual quiet
inspectable at the middle names; heartbeat choice gave the current support a
measured order of expenditure. The refining tower then funged those residual
reports across stages: each supported record could pass its admitted
information forward, and the terminal residual could be placed below a bound
that tends to zero.

That is real discipline. A nonnegative residual magnitude lies under an
accepted vanishing bound, and the apparatus can say so without pretending that
one finite support has become the whole continuum. The tower keeps the
residual report honest: it names the supports, carries the maps between them,
and records exactly where the squeeze applies. What it does not supply is the
single quantity that interprets the remaining order-two pressure. A bound may
force residual size downward while still leaving open what the residual presses
against and why one reading should organize the next stage.

The bridge out of the chapter is therefore not another support rule. Another
coordinate convention would only add names. Another local selection event would
only spend the current support again. The next demand has to gather residual
pressure into one finite quantity without erasing the supports, coordinates,
choices, and tower maps that made the pressure legible. It must preserve the
discipline Chapter 7 earned while refusing to treat a collection of local
reports as the invariant itself.

The next chapter begins from that demand. It does not abandon finite support,
and it does not smuggle in a completed analytic world behind the tower. It asks
for the invariant that can measure the remaining order-two pressure while
preserving the finite discipline just earned.

\section*{Coda: The Bolt and the Fitting}

A tailor cutting a suit from a single bolt of cloth performs this chapter's discipline with
shears and chalk: a finite field worked in hand, named before it is cut, joined by local
seams, spent by measured decisions, and squeezed toward an exact fit by successive fittings.
The trade is the chapter, in wool.

The bolt is the support, and it is finite. The tailor cuts the garment from the cloth on the
table; there is no reaching outside the bolt for more cloth in the middle of the cut. What the
bolt can yield is what the garment can be, and a tailor who has laid out the pieces knows
before cutting which cuts the cloth admits. The field of possible pieces is fixed first, the
order of cutting second --- support before order, because the cloth's capacity is settled
before the sequence of shears begins. A cut is made because the bolt can yield it, not
because the shears happened to reach that part of the table.

The pattern is the field's names. Before cutting, the tailor chalks the pattern pieces onto
the cloth --- the front, the back, the sleeve, the collar --- placing and labelling each on the
finite field. The chalk gives the bolt coordinates: where each piece sits, how adjacent pieces
face along the grain, which edges will meet. It does not pretend the bolt is endless; it names
the positions the finite cloth actually offers. The garment is addressed on the cloth before a
single thread is cut.

The seams are the local equations. A seam is not a statement about the whole garment; it is a
local condition where two named edges must match --- this edge to that edge, this dart taking
in exactly this much at exactly this point. Each seam and each dart is an equation at a place,
joining the named pieces locally, and the garment is articulate because its finite pieces
carry these local conditions rather than merely lying side by side. The cloth, named and
seamed, can carry the shape of a body the way the support, named and equationed, can carry the
shape of a residual.

The cutting is the measured selector. The tailor does not cut arbitrarily; each cut is chosen
by a carried measure --- the customer's measurements, the fall of the grain, the cloth already
spent --- spending the bolt deliberately and passing the remaining cloth forward. Where two
cuts would serve, the tailor's measure decides between them, and where the cloth is tied the
tailor leaves the choice visible for the fitting to settle. Selection spends the support by a
finite cost, exactly as the apparatus spends its support by a carried measure rather than by
whim.

And the fittings are the squeeze. No suit is exact at the first cut. The tailor bastes a rough
garment, fits it to the body, and reads the misfit --- a gap here, a pull there. Then the
garment is taken in, fitted again, and the misfit is smaller; taken in once more, and smaller
still. Each fitting is a stage in a refining tower, and the gap between the garment and the
body is driven below any tolerance the tailor cares to name, not by reaching a Platonic
perfect cloth but by a finite sequence of finite corrections. The suit is never the whole
continuum of possible garments; it is one exact garment, squeezed to fit by successive finite
fittings, with the residual misfit driven toward zero along the tower of corrections.

That is the chapter on the cutting table. The bolt is the finite support; the chalked pattern
is the field's coordinate names; the seams and darts are the local equations; the measured
cuts are the selector spending the support; and the successive fittings are the squeeze that
drives the misfit below any bound. The apparatus, like the tailor, works a finite field in
hand --- naming it, seaming it, spending it by measure, and fitting it down to an exact
residual --- and never once pretends the bolt is infinite or the fit is anything but earned.
